\section{AHSME 1999}
        \begin{enumerate}
        	\item (\textit{AHSME 1999}) $1 - 2 + 3 -4 + \cdots - 98 + 99 =$


 $\mathrm{(A) \ -50 } \qquad \mathrm{(B) \ -49 } \qquad \mathrm{(C) \ 0 } \qquad \mathrm{(D) \ 49 } \qquad \mathrm{(E) \ 50 }$


 

	\item (\textit{AHSME 1999}) Which of the following statements is false?


 $\text{(A) All equilateral triangles are congruent to each other.}$


 $\text{(B) All equilateral triangles are convex.}$


 $\text{(C) All equilateral triangles are equiangular.}$


 $\text{(D) All equilateral triangles are regular polygons.}$


 $\text{(E) All equilateral triangles are similar to each other.}$


 

	\item (\textit{AHSME 1999}) The number halfway between $1/8$ and $1/10$ is 


 $\mathrm{(A) \  } \frac 1{80} \qquad \mathrm{(B) \  } \frac 1{40} \qquad \mathrm{(C) \  } \frac 1{18} \qquad \mathrm{(D) \  } \frac 1{9} \qquad \mathrm{(E) \  } \frac 9{80}$


 

	\item (\textit{AHSME 1999}) Find the sum of all prime numbers between $1$ and $100$ that are simultaneously $1$ greater than a multiple of $4$ and $1$ less than a multiple of $5$. 


 $\mathrm{(A) \ } 118 \qquad \mathrm{(B) \ }137 \qquad \mathrm{(C) \ } 158 \qquad \mathrm{(D) \ } 187 \qquad \mathrm{(E) \ } 245$


 

	\item (\textit{AHSME 1999}) The marked price of a book was $30 \%$ less than the suggested retail price. Alice purchased the book for half the marked price at a Fiftieth Anniversary sale. What percent of the suggested retail price did Alice pay?


 $\mathrm{(A) \ }25 \% \qquad \mathrm{(B) \ }30 \% \qquad \mathrm{(C) \ }35 \% \qquad \mathrm{(D) \ }60 \% \qquad \mathrm{(E) \ }65 \%$


 

	\item (\textit{AHSME 1999}) What is the sum of the digits of the decimal form of the product $2^{1999} \cdot 5^{2001}$?


 $\mathrm{(A) \ }2 \qquad \mathrm{(B) \ }4 \qquad \mathrm{(C) \ }5 \qquad \mathrm{(D) \ }7 \qquad \mathrm{(E) \ }10$


 

	\item (\textit{AHSME 1999}) What is the largest number of acute angles that a convex hexagon can have? 


 $\mathrm{(A) \ }2 \qquad \mathrm{(B) \ }3 \qquad \mathrm{(C) \ }4 \qquad \mathrm{(D) \ }5 \qquad \mathrm{(E) \ }6$


 

	\item (\textit{AHSME 1999}) At the end of 1994 Walter was half as old as his grandmother. The sum of the years in which they were born is 3838. How old will Walter be at the end of 1999?


 $\mathrm{(A) \ } 48 \qquad \mathrm{(B) \ }49 \qquad \mathrm{(C) \ }53 \qquad \mathrm{(D) \ }55 \qquad \mathrm{(E) \ } 101$


 

	\item (\textit{AHSME 1999}) Before Ashley started a three-hour drive, her car's odometer reading was 29792, a palindrome. (A palindrome is a number that reads the same way from left to right as it does from right to left). At her destination, the odometer reading was another palindrome. If Ashley never exceeded the speed limit of 75 miles per hour, which of the following was her greatest possible average speed?


 $\mathrm{(A) \ } 33\frac 13 \qquad \mathrm{(B) \ }53\frac 13 \qquad \mathrm{(C) \ }66\frac 23 \qquad \mathrm{(D) \ }70\frac 13 \qquad \mathrm{(E) \ } 74\frac 13$


 

	\item (\textit{AHSME 1999}) A sealed envelope contains a card with a single digit on it. Three of the following statements are true, and the other is false.


 I. The digit is 1.


 II. the digit is not 2. 


 III. The digit is 3.


 IV. The digit is not 4.


 Which one of the following must necessarily be correct?


 $\mathrm{(A) \ I\ is\ true} \qquad \mathrm{(B) \ I\ is\ false} \qquad \mathrm{(C) \ II\ is\ true} \qquad \mathrm{(D) \ III\ is\ true} \qquad \mathrm{(E) \ IV\ is\ false}$


 

	\item (\textit{AHSME 1999}) The student lockers at Olympic High are numbered consecutively beginning with locker number $1$. The plastic digits used to number the lockers cost two cents apiece. Thus, it costs two cents to label locker number $9$ and four cents to label locker number $10$. If it costs $137.94$ to label all the lockers, how many lockers are there at the school?


 $\mathrm{(A) \ }2001 \qquad \mathrm{(B) \ }2010 \qquad \mathrm{(C) \ }2100 \qquad \mathrm{(D) \ }2726 \qquad \mathrm{(E) \ }6897$


 

	\item (\textit{AHSME 1999}) What is the maximum number of points of intersection of the graphs of two different fourth degree polynomial functions $y = p(x)$ and $y = q(x)$, each with leading coefficient $1$?


 $\textrm{(A)} \ 1 \qquad \textrm{(B)} \ 2 \qquad \textrm{(C)} \ 3 \qquad \textrm{(D)} \ 4 \qquad \textrm{(E)} \ 8$


 

	\item (\textit{AHSME 1999}) Define a sequence of real numbers $a_1, a_2, a_3, \ldots$ by $a_1 = 1$ and $a_{n+1}^3 = 99a_n^3$ for all $n \ge 1$. Then $a_{100}$ equals


 $\mathrm{(A) \ } 33^{33} \qquad \mathrm{(B) \ } 33^{99} \qquad \mathrm{(C) \ } 99^{33} \qquad \mathrm{(D) \ }99^{99} \qquad \mathrm{(E) \ none\ of\ the\ above}$


 

	\item (\textit{AHSME 1999}) Four girls - Mary, Aline, Tina, and Hana - sang songs in a concert as trios, with one girl sitting out each time. Hanna sang 7 songs, which was more than any other girl, and Mary sang 4 songs, which was fewer than any other girl. How many songs did these trios sing? 


 $\mathrm{(A) \ 7 } \qquad \mathrm{(B) \ 8 } \qquad \mathrm{(C) \ 9 } \qquad \mathrm{(D) \ 10 } \qquad \mathrm{(E) \ 11 }$


 

	\item (\textit{AHSME 1999}) Let $x$ be a real number such that $\sec x - \tan x = 2$. Then $\sec x + \tan x =$


 $\mathrm{(A) \ } 0.1 \qquad \mathrm{(B) \ } 0.2 \qquad \mathrm{(C) \ } 0.3 \qquad \mathrm{(D) \ } 0.4 \qquad \mathrm{(E) \ } 0.5$


 

	\item (\textit{AHSME 1999}) What is the radius of a circle inscribed in a rhombus with diagonals of length $10$ and $24$?


 $\mathrm{(A) \ }4 \qquad \mathrm{(B) \ }\frac {58}{13} \qquad \mathrm{(C) \ }\frac{60}{13} \qquad \mathrm{(D) \ }5 \qquad \mathrm{(E) \ }6$


 

	\item (\textit{AHSME 1999}) Let $P(x)$ be a polynomial such that when $P(x)$ is divided by $x-19$, the remainder is $99$, and when $P(x)$ is divided by $x - 99$, the remainder is $19$. What is the remainder when $P(x)$ is divided by $(x-19)(x-99)$?


 $\mathrm{(A) \ } -x + 80 \qquad \mathrm{(B) \ } x + 80 \qquad \mathrm{(C) \ } -x + 118 \qquad \mathrm{(D) \ } x + 118 \qquad \mathrm{(E) \ } 0$


 

	\item (\textit{AHSME 1999}) How many zeros does $f(x) = \cos(\log x)$ have on the interval $0 < x < 1$?


 $\mathrm{(A) \ } 0 \qquad \mathrm{(B) \ } 1 \qquad \mathrm{(C) \ } 2 \qquad \mathrm{(D) \ } 10 \qquad \mathrm{(E) \ } \text{infinitely\ many}$


 

	\item (\textit{AHSME 1999}) Consider all triangles $ABC$ satisfying in the following conditions: $AB = AC$, $D$ is a point on $\overline{AC}$ for which $\overline{BD} \perp \overline{AC}$, $AC$ and $CD$ are integers, and $BD^{2} = 57$.  Among all such triangles, the smallest possible value of $AC$ is


 
                    \begin{figure}[H]
                        \centering
                        \includegraphics[width=0.5\textwidth]{Problem/AHSME_I_1999_19_Problem_Figure_1.png}
                    \end{figure}
                    


 $\textrm{(A)} \ 9 \qquad \textrm{(B)} \ 10 \qquad \textrm{(C)} \ 11 \qquad \textrm{(D)} \ 12 \qquad \textrm{(E)} \ 13$


 

	\item (\textit{AHSME 1999}) The sequence $a_{1},a_{2},a_{3},\ldots$ satisfies $a_{1} = 19,a_{9} = 99$, and, for all $n\geq 3$, $a_{n}$ is the arithmetic mean of the first $n - 1$ terms.  Find $a_2$.


 $\textrm{(A)} \ 29 \qquad \textrm{(B)} \ 59 \qquad \textrm{(C)} \ 79 \qquad \textrm{(D)} \ 99 \qquad \textrm{(E)} \ 179$


 

	\item (\textit{AHSME 1999}) A circle is circumscribed about a triangle with sides $20,21,$ and $29,$ thus dividing the interior of the circle into four regions. Let $A,B,$ and $C$ be the areas of the non-triangular regions, with $C$ be the largest. Then


 $\mathrm{(A) \ }A+B=C \qquad \mathrm{(B) \ }A+B+210=C \qquad \mathrm{(C) \ }A^2+B^2=C^2 \qquad \mathrm{(D) \ }20A+21B=29C \qquad \mathrm{(E) \ } \frac 1{A^2}+\frac 1{B^2}= \frac 1{C^2}$


 

	\item (\textit{AHSME 1999}) The graphs of $y = -|x-a| + b$ and $y = |x-c| + d$ intersect at points $(2,5)$ and $(8,3)$. Find $a+c$.


 $\mathrm{(A) \ } 7 \qquad \mathrm{(B) \ } 8 \qquad \mathrm{(C) \ } 10 \qquad \mathrm{(D) \ } 13\qquad \mathrm{(E) \ } 18$


 

	\item (\textit{AHSME 1999}) The equiangular convex hexagon $ABCDEF$ has $AB = 1, BC = 4, CD = 2,$ and $DE = 4.$ The area of the hexagon is 


 $\mathrm{(A) \ } \frac {15}2\sqrt{3} \qquad \mathrm{(B) \ }9\sqrt{3} \qquad \mathrm{(C) \ }16 \qquad \mathrm{(D) \ }\frac{39}4\sqrt{3} \qquad \mathrm{(E) \ } \frac{43}4\sqrt{3}$


 

	\item (\textit{AHSME 1999}) Six points on a circle are given. Four of the chords joining pairs of the six points are selected at random. What is the probability that the four chords form a convex quadrilateral?


 $\mathrm{(A) \ } \frac 1{15} \qquad \mathrm{(B) \ } \frac 1{91} \qquad \mathrm{(C) \ } \frac 1{273} \qquad \mathrm{(D) \ } \frac 1{455} \qquad \mathrm{(E) \ } \frac 1{1365}$


 

	\item (\textit{AHSME 1999}) There are unique integers $a_{2},a_{3},a_{4},a_{5},a_{6},a_{7}$ such that


 
 \[\frac {5}{7} = \frac {a_{2}}{2!} + \frac {a_{3}}{3!} + \frac {a_{4}}{4!} + \frac {a_{5}}{5!} + \frac {a_{6}}{6!} + \frac {a_{7}}{7!}\]


 where $0\leq a_{i} < i$ for $i = 2,3,\ldots,7$.  Find $a_{2} + a_{3} + a_{4} + a_{5} + a_{6} + a_{7}$.


 $\textrm{(A)} \ 8 \qquad \textrm{(B)} \ 9 \qquad \textrm{(C)} \ 10 \qquad \textrm{(D)} \ 11 \qquad \textrm{(E)} \ 12$


 

	\item (\textit{AHSME 1999}) Three non-overlapping regular plane polygons, at least two of which are congruent, all have sides of length $1$. The polygons meet at a point $A$ in such a way that the sum of the three interior angles at $A$ is $360^{\circ}$. Thus the three polygons form a new polygon with $A$ as an interior point. What is the largest possible perimeter that this polygon can have? 


 $\mathrm{(A) \ }12 \qquad \mathrm{(B) \ }14 \qquad \mathrm{(C) \ }18 \qquad \mathrm{(D) \ }21 \qquad \mathrm{(E) \ } 24$


 

	\item (\textit{AHSME 1999}) In triangle $ABC$, $3 \sin A + 4 \cos B = 6$ and $4 \sin B + 3 \cos A = 1$. Then $\angle C$ in degrees is 


 $\mathrm{(A) \ }30 \qquad \mathrm{(B) \ }60 \qquad \mathrm{(C) \ }90 \qquad \mathrm{(D) \ }120 \qquad \mathrm{(E) \ }150$


 

	\item (\textit{AHSME 1999}) Let $x_1, x_2, \ldots , x_n$ be a sequence of integers such that


 $\text{(i)}$ $-1 \le x_i \le 2$ $\text{for}$ $i = 1,2, \ldots n$


 $\text{(ii)}$ $x_1 + \cdots + x_n = 19$; $\text{and}$


 $\text{(iii)}$ $x_1^2 + x_2^2 + \cdots + x_n^2 = 99$.


 Let $m$ and $M$ be the minimal and maximal possible values of $x_1^3 + \cdots + x_n^3$, respectively. Then $\frac Mm =$


 $\mathrm{(A) \ }3 \qquad \mathrm{(B) \ }4 \qquad \mathrm{(C) \ }5 \qquad \mathrm{(D) \ }6 \qquad \mathrm{(E) \ }7$


 

	\item (\textit{AHSME 1999}) A tetrahedron with four equilateral triangular faces has a sphere inscribed within it and a sphere circumscribed about it. For each of the four faces, there is a sphere tangent externally to the face at its center and to the circumscribed sphere. A point $P$ is selected at random inside the circumscribed sphere. The probability that $P$ lies inside one of the five small spheres is closest to 


 $\mathrm{(A) \ }0 \qquad \mathrm{(B) \ }0.1 \qquad \mathrm{(C) \ }0.2 \qquad \mathrm{(D) \ }0.3 \qquad \mathrm{(E) \ }0.4$


 

	\item (\textit{AHSME 1999}) The number of ordered pairs of integers $(m,n)$ for which $mn \ge 0$ and


 
 \[m^3 + n^3 + 99mn = 33^3\]


 is equal to


 $\mathrm{(A) \ }2 \qquad \mathrm{(B) \ } 3\qquad \mathrm{(C) \ } 33\qquad \mathrm{(D) \ }35 \qquad \mathrm{(E) \ } 99$


 

\end{enumerate}